\documentclass[11pt]{article}
\usepackage[margin=25mm]{geometry}
\usepackage{fontspec}\setmainfont{Latin Modern Roman}
\usepackage{amsmath,amssymb,mathtools}
\usepackage{unicode-math}\setmathfont{Latin Modern Math}
\usepackage{xurl,hyperref}
\hypersetup{colorlinks=true,urlcolor=blue}
\setcounter{secnumdepth}{0}
\setlength{\parindent}{0pt}\setlength{\parskip}{5pt}
\setlength{\emergencystretch}{3em}
\newcommand{\sref}[2]{\hyperlink{source:#1:#2}{[S#2]}}
\newcommand{\cataloguescope}[1]{\paragraph{Recorded target.}#1}
\begin{document}
\section{Kourovka Notebook Problem 21.126}\label{statement}

\subsection{Definitions and mathematical statement}

\paragraph{Shared notation.}$\mathbb N=\{1,2,\ldots\}$, and $\mathbb Z,\mathbb Q,\mathbb R,\mathbb C$ denote the integers, rationals, reals and complex numbers. Any other conventions are specified locally.


For a finite graph $\Lambda$, its right-angled Artin group has generators $v\in V(\Lambda)$ and relations $[v,w]=1$ for the edges. For a finite presentation of a group $H$, its Dehn function $\delta_H(n)$ is the greatest minimal number of relators needed to fill a null word of length at most $n$. Functions are compared up to the usual equivalence generated by $f\preccurlyeq g$ when $f(n)\leq Cg(Cn+C)+Cn+C$ for some $C>0$. Here super-exponential means faster than every exponential, and sub-exponential but nonpolynomial means exponential rate zero while no polynomial upper bound exists, in this coarse sense.
\paragraph{Statement recorded in the supplied source.}
Do finitely presented subgroups $H$ of right-angled Artin groups exist whose Dehn functions are (a) super-exponential, or (b) sub-exponential and not polynomially bounded? \sref{statement}{1}\sref{statement}{2}
\subsection{Short English statement}
Find finitely presented subgroups of right-angled Artin groups with either faster-than-exponential filling or slower-than-exponential filling that still exceeds every polynomial.
\subsection{Sources}
\begin{description}
\item[\hypertarget{source:statement:1}{S1}] ULAM.AI and the UnsolvedMath Contributors, \emph{UnsolvedMath: A Curated Collection of Open Mathematics Problems}, catalogue entry. CC BY 4.0. \url{https://huggingface.co/datasets/ulamai/UnsolvedMath}.
\item[\hypertarget{source:statement:2}{S2}] E. I. Khukhro and V. D. Mazurov (editors), The Kourovka Notebook, Unsolved Problems in Group Theory, No.21, version48 (6 October2026), Problem 21.126 and its remarks. \url{https://arxiv.org/abs/1401.0300v48}.
\end{description}
\label{end:statement}
\end{document}
